% Generated from validated Article Draft Result. Do not hand-edit; revise the structured draft instead.
\section{Detailed Chronology――圧縮した物語の外側に、出来事を残す}
\noindent\textbf{年間本文はstoryとtrajectoryへ圧縮するが、歴史記録としてはrelease、paper initial publication、preview、API/product availability、policy/control milestoneを同一視できない。本節では、年次叙述で失われやすいevent-levelの差を読むための基準を明示する。}

年間の出来事は、研究論文の初出、モデル発表、preview、一般提供、API利用可能化、製品統合、policy公開など複数の状態に分かれる。物語上ひとつの系列にまとめる場合でも、これらを同じ日付・同じ公開状態へ潰してはならない。


後年に広く普及した機能であっても、2023年のある時点で同じ範囲に提供されていたとは限らない。年末時点の状態や2026年から見た重要性を、年初の利用可能性へ逆投影しないことをchronologyの基本原則とする。


一次資料が日付まで示す場合と、月単位・publication lifecycleのみが確認できる場合を区別する。精度が不足する出来事に推測で日付を補わず、sourceが持つ時間精度をそのまま残す。


本節のevent indexは、受理済みScreeningで明示的なpublished\_atを持つretain済み記録から決定論的に再展開する。日付不明の記録は推測して補わず監査に未解決として残し、source locatorは読者向け本文とは別のauditで保持する。
